\chapter{Cài đặt bằng OpenMP}
\label{ch:impl}

Chương này mô tả cách các thuật toán ở \cref{ch:algorithms} được hiện thực bằng C11 và OpenMP:
cấu trúc mã nguồn (\cref{sec:impl-arch}), chi tiết từng phiên bản (\cref{sec:impl-versions}),
các kỹ thuật tối ưu (\cref{sec:impl-opt}), kiểm thử (\cref{sec:impl-test}) và công cụ đo hiệu
năng (\cref{sec:impl-bench}). Toàn bộ mã nguồn có trong \cref{app:code} và trên kho GitHub của
nhóm.

\section{Kiến trúc mã nguồn}
\label{sec:impl-arch}

\begin{table}[H]
\centering
\caption{Các tệp mã nguồn.}
\label{tab:files}
\renewcommand{\arraystretch}{1.2}
\begin{tabularx}{\textwidth}{>{\ttfamily}p{4cm}X}
\toprule
\textnormal{\textbf{Tệp}} & \textbf{Nội dung} \\
\midrule
src/scan.h & Giao diện chung: kiểu \code{elem\_t}, cấu trúc vùng làm việc \code{scan\_ws},
nguyên mẫu các hàm scan, bảng đăng ký thuật toán. \\
src/scan.c & Cài đặt sáu phiên bản scan và kernel \code{copy} tham chiếu. \\
src/test\_scan.c & Bộ kiểm thử tính đúng: so sánh mọi phiên bản với bản tuần tự. \\
src/bench.c & Chương trình đo hiệu năng, xuất CSV. \\
src/compact.c & Ứng dụng stream compaction dựa trên scan. \\
scripts/run\_all.sh & Chạy toàn bộ thực nghiệm, ghi thông tin máy và kết quả thô. \\
scripts/plots.py & Tổng hợp trung vị, vẽ biểu đồ và sinh bảng LaTeX. \\
Makefile & Biên dịch với \code{gcc -O3 -march=native -fopenmp}. \\
\bottomrule
\end{tabularx}
\end{table}

\paragraph{Giao diện thống nhất.} Mọi phiên bản có cùng chữ ký (\cref{lst:iface}), cho phép
chương trình đo và bộ kiểm thử duyệt qua một bảng các con trỏ hàm mà không cần mã riêng cho từng
thuật toán. Bộ nhớ phụ (buffer cho Hillis--Steele và Blelloch) được cấp phát \emph{trước} khi đo
qua \code{scan\_ws\_create}, để thời gian đo không bao gồm chi phí \code{malloc} và lỗi trang.

\begin{lstlisting}[caption={Giao diện chung trong \texttt{scan.h} (rút gọn).},label={lst:iface}]
typedef int64_t elem_t;
typedef struct {
    elem_t *tmp;      // buffer phụ >= next_pow2(n) phần tử
    size_t  tmp_cap;
    size_t  block;    // kích thước khối B cho blocked_opt
} scan_ws;

typedef void (*scan_fn)(const elem_t *in, elem_t *out, size_t n, scan_ws *ws);

void scan_seq          (const elem_t *in, elem_t *out, size_t n, scan_ws *ws);
void scan_hillis_steele(const elem_t *in, elem_t *out, size_t n, scan_ws *ws);
void scan_blelloch     (const elem_t *in, elem_t *out, size_t n, scan_ws *ws);
void scan_blocked      (const elem_t *in, elem_t *out, size_t n, scan_ws *ws);
void scan_blocked_opt  (const elem_t *in, elem_t *out, size_t n, scan_ws *ws);
void scan_omp          (const elem_t *in, elem_t *out, size_t n, scan_ws *ws);
\end{lstlisting}

\paragraph{Cấp phát bộ nhớ.} Mọi mảng được cấp phát bằng \code{aligned\_alloc(64, \dots)}, căn
lề theo dòng cache, để các khối của luồng bắt đầu đúng biên dòng cache và trình biên dịch có thể
dùng lệnh nạp vector căn lề.

\section{Cài đặt các phiên bản}
\label{sec:impl-versions}

\subsection{Hillis--Steele}

Cài đặt (\cref{lst:hs}) mở \emph{một} vùng \code{parallel} duy nhất bao trọn mọi bước, tránh chi
phí fork--join \(\log_2 n\) lần. Mỗi bước tách làm hai vòng \code{for}: vòng \(i < \mathit{off}\)
chỉ sao chép, vòng \(i \ge \mathit{off}\) thực hiện phép cộng. Việc tách bỏ nhánh \code{if} trong
thân vòng lặp giúp trình biên dịch vector hóa được. Vòng thứ nhất dùng \code{nowait} vì hai vòng
ghi vào hai vùng rời nhau của \code{b}. Việc hoán đổi con trỏ đặt trong \code{single}: mọi luồng
thấy cùng giá trị \code{a}, \code{b} vì đây là biến dùng chung khai báo ngoài vùng song song, và
rào chắn ngầm định cuối \code{single} đảm bảo không luồng nào bắt đầu bước sau khi hoán đổi chưa
xong.

\begin{lstlisting}[caption={Hillis--Steele với bộ đệm kép.},label={lst:hs}]
elem_t *a = out, *b = ws->tmp;
#pragma omp parallel
{
    #pragma omp for schedule(static)
    for (size_t i = 0; i < n; i++) a[i] = in[i];

    for (size_t off = 1; off < n; off <<= 1) {
        #pragma omp for schedule(static) nowait
        for (size_t i = 0; i < off; i++) b[i] = a[i];
        #pragma omp for schedule(static)
        for (size_t i = off; i < n; i++) b[i] = a[i - off] + a[i];
        #pragma omp single
        { elem_t *t = a; a = b; b = t; }      // barrier ngầm định
    }
    if (a != out) {
        #pragma omp for schedule(static)
        for (size_t i = 0; i < n; i++) out[i] = a[i];
    }
}
\end{lstlisting}

\subsection{Blelloch}

Phiên bản Blelloch (\cref{lst:blelloch}) làm việc trên buffer \code{ws->tmp} có kích thước
\(m = 2^{\lceil\log_2 n\rceil}\); phần đệm \(x[n..m-1]\) được gán 0 (phần tử đơn vị) nên không
ảnh hưởng kết quả. Vòng lặp theo \code{k} dùng bước nhảy \code{stride} thay vì nhân chỉ số, để số
lần lặp của mỗi tầng là \(m/\mathit{stride}\) và \code{schedule(static)} chia đều chúng. Ở các
tầng trên cùng, số lần lặp ít hơn số luồng, nên phần lớn luồng không có việc nhưng vẫn phải chờ
ở rào chắn ngầm định cuối mỗi vòng \code{for}. Bước cuối cộng \code{in[i]} để chuyển exclusive
thành inclusive.

\begin{lstlisting}[caption={Blelloch: up-sweep và down-sweep tại chỗ.},label={lst:blelloch}]
const size_t m = next_pow2(n);
elem_t *x = ws->tmp;
#pragma omp parallel
{
    #pragma omp for schedule(static)
    for (size_t i = 0; i < m; i++) x[i] = (i < n) ? in[i] : 0;

    for (size_t stride = 2; stride <= m; stride <<= 1) {      // up-sweep
        const size_t half = stride >> 1;
        #pragma omp for schedule(static)
        for (size_t k = stride - 1; k < m; k += stride)
            x[k] += x[k - half];
    }
    #pragma omp single
    x[m - 1] = 0;
    for (size_t stride = m; stride >= 2; stride >>= 1) {      // down-sweep
        const size_t half = stride >> 1;
        #pragma omp for schedule(static)
        for (size_t k = stride - 1; k < m; k += stride) {
            elem_t t    = x[k - half];
            x[k - half] = x[k];
            x[k]       += t;
        }
    }
    #pragma omp for schedule(static)
    for (size_t i = 0; i < n; i++) out[i] = x[i] + in[i];    // -> inclusive
}
\end{lstlisting}

\subsection{Blocked}

Phiên bản chia khối cơ bản (\cref{lst:blocked}) không dùng \omp{for} mà tự tính khoảng
\([lo, hi)\) của mỗi luồng từ \code{omp\_get\_thread\_num()}. Lý do là pha 1 và pha 3 phải dùng
\emph{cùng một} cách chia; tự chia giúp điều này luôn đúng, không phụ thuộc vào cách thư viện cài
đặt \code{schedule(static)}. Mảng \code{sums} có \(p + 1\) phần tử, \code{sums[0] = 0}, nên sau khi
scan tuần tự thì \code{sums[t]} chính là offset của luồng \(t\). Thuật toán chỉ cần hai điểm đồng
bộ: \code{barrier} sau pha 1 và rào chắn ngầm định cuối \code{single}.

\begin{lstlisting}[caption={Blocked: ba pha, mỗi luồng một khối liên tiếp.},label={lst:blocked}]
elem_t sums[MAX_THREADS + 1];
#pragma omp parallel
{
    const int    p   = omp_get_num_threads();
    const int    t   = omp_get_thread_num();
    const size_t len = (n + p - 1) / p;
    const size_t lo  = min(t * len, n), hi = min(lo + len, n);

    elem_t s = 0;                                   // Pha 1
    for (size_t i = lo; i < hi; i++) s += in[i];
    sums[t + 1] = s;
    #pragma omp barrier

    #pragma omp single                              // Pha 2
    { sums[0] = 0;
      for (int i = 1; i <= p; i++) sums[i] += sums[i - 1]; }

    s = sums[t];                                    // Pha 3
    for (size_t i = lo; i < hi; i++) { s += in[i]; out[i] = s; }
}
\end{lstlisting}

\subsection{Blocked-opt}

Phiên bản tối ưu (\cref{lst:blockedopt}) hiện thực ý tưởng chia tile ở \cref{sec:alg-tiled}.
Ba điểm khác so với Blocked:
\begin{enumerate}
  \item \textbf{Vòng lặp tile:} mỗi luồng xử lý khối thứ \(t\) của mỗi tile; mỗi khối \(B\)
        phần tử (mặc định \(B = 2^{15}\), tức 256~KiB, vừa L2 của cả P-core lẫn E-core).
  \item \textbf{Không có \code{single}:} thay vì một luồng scan mảng tổng khối rồi phát lại cho
        các luồng khác, \emph{mỗi luồng tự cộng} các tổng khối đứng trước nó (\(\Oh(p)\) phép
        cộng, \(p \le 12\)) và tự cập nhật bản sao \code{carry} riêng. Nhờ đó mỗi tile chỉ cần
        \textbf{một} rào chắn.
  \item \textbf{Bộ đệm kép cho tổng khối:} tile chẵn ghi vào \code{part[0]}, tile lẻ ghi vào
        \code{part[1]}. Một luồng nhanh đã sang tile kế tiếp sẽ ghi vào buffer kia, không ghi đè
        dữ liệu mà luồng chậm còn đang đọc. Đây là lý do có thể bỏ rào chắn thứ hai.
\end{enumerate}

\begin{lstlisting}[caption={Blocked-opt: tile vừa L2, một barrier mỗi tile.},label={lst:blockedopt}]
typedef struct { elem_t v; char pad[64 - sizeof(elem_t)]; } padded_elem;
static padded_elem part[2][MAX_THREADS];          // đệm 64 byte, bộ đệm kép

#pragma omp parallel
{
    const int p = omp_get_num_threads(), t = omp_get_thread_num();
    const size_t tile = B * p;
    elem_t carry = 0;  int buf = 0;
    for (size_t base = 0; base < n; base += tile) {
        const size_t lo = min(base + t * B, n), hi = min(lo + B, n);
        elem_t s = 0;
        #pragma omp simd reduction(+:s)             // Pha 1 (AVX2)
        for (size_t i = lo; i < hi; i++) s += in[i];
        part[buf][t].v = s;
        #pragma omp barrier
        elem_t off = carry, tot = 0;                // Pha 2: mỗi luồng tự tính
        for (int j = 0; j < p; j++) {
            elem_t v = part[buf][j].v;
            if (j < t) off += v;
            tot += v;
        }
        carry += tot;
        s = off;                                    // Pha 3: dữ liệu còn trong L2
        for (size_t i = lo; i < hi; i++) { s += in[i]; out[i] = s; }
        buf ^= 1;
    }
}
\end{lstlisting}

\paragraph{Tính đúng của việc bỏ rào chắn thứ hai.} Xét luồng \(u\) đã xong tile \(k\) và bắt
đầu tile \(k+1\), ghi vào \code{part[(k+1) mod 2]}. Buffer này không bị đọc ở tile \(k\). Luồng
\(u\) chỉ có thể sang tile \(k+2\), tức ghi lại vào \code{part[k mod 2]}, sau khi vượt qua rào
chắn của tile \(k+1\). Mà rào chắn đó chỉ mở khi \emph{mọi} luồng đã tới, tức mọi luồng đã đọc
xong \code{part[k mod 2]} ở pha 2 của tile \(k\). Vậy không có tranh chấp dữ liệu.

\subsection{Chỉ thị \texttt{omp scan}}

Phiên bản này chỉ gồm vài dòng như \cref{lst:omp-scan-intro}. Trong quá trình cài đặt, nhóm gặp
hai vấn đề với GCC 15.2:
\begin{enumerate}
  \item Ban đầu nhóm dùng \code{parallel for simd reduction(inscan,+:s)} để kết hợp đa luồng và
        vector hóa. Với \(n\) lớn, chương trình bị \textbf{tràn stack}: AddressSanitizer chỉ ra
        GCC cấp phát một mảng tạm kích thước \(\Oh(n)\) trên stack cho phần \code{simd}. Nhóm bỏ
        \code{simd} để khắc phục.
  \item Khi \(n = 0\) và \(p > 1\), libgomp \textbf{ghi đè bộ nhớ} (lỗi
        ``\emph{stack smashing detected}''). Bộ kiểm thử phát hiện lỗi này; nhóm thêm kiểm tra
        \code{if (n == 0) return;} ở đầu hàm.
\end{enumerate}
Hai vấn đề này cho thấy chỉ thị \code{scan} tuy tiện lợi nhưng còn khá mới, và cài đặt trong
trình biên dịch phổ biến vẫn có lỗi ở trường hợp biên. Đây cũng là lý do bộ kiểm thử phải bao
gồm các kích thước đặc biệt như 0 và 1.

\section{Các kỹ thuật tối ưu}
\label{sec:impl-opt}

\begin{table}[H]
\centering
\caption{Tổng hợp các kỹ thuật tối ưu đã áp dụng.}
\label{tab:opts}
\renewcommand{\arraystretch}{1.25}
\begin{tabularx}{\textwidth}{>{\raggedright\arraybackslash}p{3.6cm}X>{\raggedright\arraybackslash}p{2.6cm}}
\toprule
\textbf{Kỹ thuật} & \textbf{Mục đích} & \textbf{Áp dụng cho} \\
\midrule
Một vùng \code{parallel} duy nhất & Tránh fork--join lặp lại ở mỗi bước/tầng. & Mọi phiên bản \\
Tách vòng lặp bỏ nhánh & Cho phép trình biên dịch vector hóa (AVX2). & Hillis--Steele \\
\code{omp simd reduction} & Vector hóa pha tính tổng khối: 4 phần tử int64 mỗi lệnh. &
Blocked-opt \\
Chia tile vừa L2 & Pha 3 đọc lại dữ liệu từ cache: lưu lượng DRAM từ \(24n\) xuống \(16n\) byte.
& Blocked-opt \\
Đệm 64 byte (padding) & Chống false sharing trên mảng tổng khối. & Blocked-opt \\
Bộ đệm kép tổng khối & Bỏ một rào chắn mỗi tile. & Blocked-opt \\
Căn lề 64 byte & Khối bắt đầu ở biên dòng cache; nạp vector căn lề. & Mọi mảng \\
First-touch song song & Trang nhớ được ánh xạ trước khi đo; không tính lỗi trang. & Mọi mảng \\
Cấp phát trước buffer & Không tính \code{malloc} vào thời gian đo. & Hillis--Steele, Blelloch \\
\code{-O3 -march=native} & Tối ưu mạnh và dùng tập lệnh của CPU đang chạy (AVX2). & Toàn bộ \\
\bottomrule
\end{tabularx}
\end{table}

\paragraph{Vì sao không vector hóa được pha 3?} Vòng \code{s += in[i]; out[i] = s;} có phụ thuộc
mang qua vòng lặp nên trình biên dịch không tự vector hóa. Về lý thuyết có thể dùng scan trong
thanh ghi SIMD bằng các lệnh dịch và cộng (in-register scan), nhưng vì scan đã bị giới hạn bởi
băng thông bộ nhớ (\cref{sec:roofline}), lợi ích không đáng kể. Thực nghiệm xác nhận bản tuần tự
vô hướng đã đạt khoảng 13~GB/s, gần bằng kernel copy đơn luồng.

\section{Kiểm thử tính đúng}
\label{sec:impl-test}

Chương trình \code{test\_scan} so sánh từng phần tử kết quả của mọi phiên bản với
\code{scan\_seq} trên một lưới tham số:
\begin{itemize}
  \item \textbf{Kích thước} \(n \in \{0, 1, 2, 3, 5, 8, 13, 64, 100, 1000, 4096, 4099, 65537,
        1000003, 2^{20}+7\}\): gồm trường hợp rỗng, một phần tử, lũy thừa của 2, số nguyên tố và
        số lớn hơn lũy thừa của 2 một chút (buộc Blelloch phải đệm gần gấp đôi).
  \item \textbf{Số luồng} \(p \in \{1, 2, 3, 4, 7, 8, 12\}\): gồm số lẻ và số không chia hết
        \(n\), để kiểm tra khối cuối ngắn hơn hoặc rỗng.
  \item \textbf{Kích thước khối} \(B \in \{1, 3, 1024, 32768\}\) cho Blocked-opt: \(B = 1\) tạo
        rất nhiều tile, kiểm tra kỹ cơ chế bộ đệm kép và carry.
  \item \textbf{Dữ liệu} ngẫu nhiên trong \([-1000, 1000]\), có cả số âm.
  \item Kiểm tra riêng ví dụ của Blelloch: \([3,1,7,0,4,1,6,3] \mapsto [3,4,11,11,15,16,22,25]\).
\end{itemize}
Tổng cộng \textbf{945 trường hợp kiểm thử, tất cả đều đạt}. Ngoài ra, mã được biên dịch với
\code{-fsanitize=address} để phát hiện truy cập ngoài vùng nhớ; cách này đã giúp tìm ra hai lỗi
của \code{omp scan} kể trên. Chương trình đo hiệu năng cũng tự kiểm tra kết quả sau lần chạy làm
nóng và dừng ngay nếu sai.

\section{Công cụ và quy trình đo hiệu năng}
\label{sec:impl-bench}

Chương trình \code{bench} nhận danh sách thuật toán, kích thước, số luồng, số lần lặp và kích
thước khối, rồi in mỗi lần đo thành một dòng CSV \code{algo,n,threads,block,rep,time\_s}. Với mỗi
cấu hình, quy trình như sau:
\begin{enumerate}
  \item Cấp phát và khởi tạo song song các mảng (first-touch), tính kết quả tham chiếu.
  \item Đặt số luồng bằng \code{omp\_set\_num\_threads(p)}.
  \item Chạy \textbf{2 lần làm nóng} (warm-up): đánh thức nhóm luồng, nạp cache lệnh, đưa CPU lên
        xung nhịp cao. Kiểm tra kết quả.
  \item Chạy \textbf{10 lần đo}, mỗi lần bao quanh bởi \code{omp\_get\_wtime()}.
\end{enumerate}
Script \code{run\_all.sh} gọi \code{bench} cho sáu nhóm thực nghiệm, ghi lại thông tin máy
(\code{lscpu}, phiên bản GCC, cờ biên dịch) cùng kết quả. Script \code{plots.py} lấy
\textbf{trung vị} của 10 lần đo cho mỗi cấu hình (ít nhạy với giá trị ngoại lai hơn trung bình,
vốn hay xuất hiện trên laptop do hệ điều hành và điều chỉnh xung nhịp), tính speedup theo
\eqref{eq:speedup}, băng thông hiệu dụng \(16n/T\) và vẽ biểu đồ. Nhờ quy trình tự động, toàn bộ
\cref{ch:experiments} có thể được tái tạo bằng hai lệnh:
\begin{lstlisting}[language=bash,numbers=none]
bash scripts/run_all.sh      # khoảng 20 phút
python scripts/plots.py      # bảng + biểu đồ vào report/
\end{lstlisting}
